%! TEX program = xelatex
\documentclass[a4paper, 11pt]{report} 

\usepackage{/home/aglaia/Documents/latex-setup/packages/packs}
\usepackage{/home/aglaia/Documents/latex-setup/packages/macros}
\usepackage{/home/aglaia/Documents/latex-setup/packages/mainstyle} 
\usepackage{/home/aglaia/Documents/latex-setup/packages/colors}    
\usepackage{/home/aglaia/Documents/latex-setup/packages/frameboxes} 
\usepackage{/home/aglaia/Documents/latex-setup/packages/title}

\usepackage{etoolbox}
\usepackage{algorithm}
\usepackage{algpseudocode}

\algrenewcommand\algorithmicwhile{\textcolor{blue}{while}}
\algrenewcommand\algorithmicdo{\textcolor{blue}{do}}
\algrenewcommand\algorithmicend{\textcolor{blue}{end}}
\algrenewcommand\algorithmicif{\textcolor{blue}{if}}
\algrenewcommand\algorithmicthen{\textcolor{blue}{then}}
\algrenewcommand\algorithmicfor{\textcolor{blue}{for}}

\setuniversity{Technische Universität Wien}
\setcoursename{Computer-Aided Verification}
\setcoursebook{slides}
\setauthorname{aglaia norza}
\setauthoremail{norzaglaia@gmail.com}
\setauthorgithub{AglaiaNorza}

\setcoverphoto{cover-meme}

\setcoverwidth{0.5\textwidth}
\setcoverheight{0.6\textwidth}
\setbeforespace{0em}
\setafterspace{0em}
\makeatletter
\patchcmd{\makecustomcover}
  {\end{tikzpicture} \par}
{\end{tikzpicture} \par \vspace{-0.5em}{\small\textcolor{chaptergray}{meme translation: ``an engineer's carreer: 1) uni life 2) first steps in the job market 3) [whatever that is] 4) surface-to-air missile builder for Leonardo S.p.A (famous italian defense company with very dubious ethics)''}} \par}
  {}{}
\makeatother

\begin{document}

\makecustomcover

\tableofcontents
\clearpage

\chapter{Introduction}

\section{Verification \& Validation}
Software and integrated circuits are pervasive, and a huge effort is spent on verification and validation (V\&V). Approximately 50\% of software development time and 35\% of hardware development time is dedicated to these tasks. The core goal of Computer-Aided Verification (CAV) is either to formally establish the correctness of a system or to find bugs.

\section{Model Checking Basics}

\begin{defframe}{Transition Systems}{}
Programs and circuits are modeled as \textbf{transition systems}.
\begin{itemize}
    \item \textbf{States ($S$):} represent configurations, such as program counter values ($pc$) and variable assignments.
    \item \textbf{Transitions ($T$):} define how states change, essentially capturing the operational semantics of the system by mapping current states to next states ($s \to s'$).
\end{itemize}
\end{defframe}

A formal model checking problem can consist of:
\begin{itemize}
    \item starting states ($i$): the initial configurations of the system;
    \item transitions ($t$): the allowable state updates;
    \item bad states ($\neg p$): states that violate the specification;
\end{itemize}

Violations of specifications are witnessed by a \textbf{counterexample}.


\begin{gframe}{Safety vs Liveness violations}
\begin{itemize}
    \item \textbf{Safety violation:} a counterexample consisting of a path of \textit{finite length} that reaches a bad state.
    \begin{itemize}
        \item example: "two intersecting traffic lights must never be green at the same time"
        \subitem counterexample: a finite path where the North-South light turns green, a timer glitches, and the East-West light also turns green - the path ends the moment this bad state is reached
    \end{itemize}
    \item \textbf{Liveness violation:} a counterexample consisting of a path of \textit{infinite length} that repeatedly visits bad states and never reaches a good state.
    \begin{itemize}
        \item example: "if a pedestrian presses the crosswalk button, the walk signal will eventually turn green"
        \subitem counterexample: an infinite loop where the pedestrian presses the button, but the traffic lights never turn green
    \end{itemize}
\end{itemize}
\end{gframe}

To properly define "bad" or "good" states, engineers verify systems against formal specifications written in \textbf{Temporal Logic}. This replaces natural language requirements with precise mathematical constraints (e.g., $GFp \to G(Gq \vee Fr)$).

\section{The State-Space Explosion Problem}
The main challenge in model checking is the exponential growth of possible states.

\begin{gframe}{Sources of state explosion}
\begin{itemize}
    \item \textbf{Data:} enumerating large domains of values, including user inputs;
    \item \textbf{Concurrency:} enumerating all possible process and thread interleaving schedules;
    \item \textbf{Environment:} enumerating all possible states of the execution environment;
    \item \textbf{Faults:} faulty processes that may feed others with junk data;
\end{itemize}
\end{gframe}

\subsection{Remedies}
To mitigate state-space explosion, several techniques are utilized:
\begin{itemize}
    \item \textbf{Abstraction:} mapping concrete values into a coarser, smaller domain;
    \item \textbf{Symbolic techniques:} using Binary Decision Diagrams (BDDs), boolean satisfiability solvers (SAT), and satisfiability-modulo-theory (SMT) solvers to handle sets of states;
    \item \textbf{Partial order reduction:} avoiding the exploration of states via "redundant paths";
    \item \textbf{Symmetry reduction:} avoiding the exploration of symmetric states;
\end{itemize}

\section{Symbolic Model Checking}
Instead of exploring states one by one (explicit state model checking), symbolic model checking processes sets of states simultaneously.

\begin{itemize}
    \item \textbf{Sets of States} are epresented as \textit{logical formulae} $F(V)$ over program variables, registers, and signals
    \item \textbf{Transitions} are represented as \textit{binary relations} $T(V, V')$ that link current variables $V$ to next-state variables $V'$
\end{itemize}

\begin{gframe}{example: reachability}
Reachability can be mathematically expressed to check if a property $P$ can be violated in $k$ steps: 
$$I(V_0) \wedge \left(\bigwedge_{i=1}^k T(V_{i-1}, V_i)\right) \wedge \neg P(V_k)$$
where
\begin{itemize}
    \item $I(V_{0})$ means that $V_{0}$ is an initial state
    \item $\bigwedge_{i=1}^k T(V_{i-1}, V_i)$ represents the (valid) transitions
    \item $\neg P(V_k)$
\end{itemize}
\end{gframe}

\subsection{Invariants and Fixpoints}
When computing reachable states iteratively, we define $R_0 = I$, and subsequent reachable sets as $T(R_{i-1}) = R_i$. The set of all states reachable in up to $k$ steps is $R_{\le k} = \bigcup_{i=0}^k R_i$.

\begin{defframe}{Safety condition}{}
A system is considered safe if:
\begin{enumerate}
    \item $R_{\le k}$ contains the initial states $I$.
    \item $T$ cannot leave $R_{\le k}$ (the system has reached a \textbf{fixpoint}).
    \item $R_{\le k}$ does not overlap with the bad states $\neg P$.
\end{enumerate}
\end{defframe}

Since exact reachability computations retard convergence and require complex quantifier alternation ($\forall\exists$), practical approaches often compute over-approximate inductive frames to detect fixpoints more efficiently.

\section{Abstraction and CEGAR}
Abstraction reduces the state space but introduces non-determinism and may introduce divergent behaviors (non-termination). An abstract program must strictly \textbf{over-approximate} the original program, meaning it contains "more" behaviors than the concrete system.

To handle this automatically, modern model checkers use the \textbf{Counterexample-Guided Abstraction Refinement (CEGAR)} loop:
\begin{enumerate}
    \item \textbf{check abstraction:} see if the abstract model is safe; if safe, the concrete system is safe
    \item \textbf{check feasibility:} if a counterexample is found, check if it is feasible in the concrete system; if feasible, a real bug is found
    \item \textbf{refine:} if the counterexample is infeasible (a false positive caused by over-approximation), use the failure trace to refine the abstraction and check again
\end{enumerate}

\section{Protocols \& Execution Models}

Verification models differ based on whether they target hardware or software.

\begin{itemize}
    \item \textbf{Concurrent Protocols (software):} standard assumptions include that processes run \textbf{asynchronously} (interleave), the scheduler is completely non-deterministic, and individual statements execute atomically; examples include threads accessing shared variables (requiring mutual exclusion algorithms like Peterson's)
    \item \textbf{Hardware Protocols:} standard assumptions include that components run \textbf{synchronously} based on a global clock (every component makes a step at a time) and communicate via actions or shared memory; examples include cache coherence protocols like MSI (Modified-Shared-Invalid)
\end{itemize}

\chapter{Temporal Logics and Model Checking}

\begin{defframe}{Finite-State Transition System}{}
    A triple $(S, T, I)$ is a \textbf{Finite-State Transition System} if 
    \begin{itemize}
        \item $S$ is a finite set of states
        \item $I \subseteq S$ is a set of initial states
        \item $T \subseteq S \times S$ is a \textit{total} transition relation ($\forall s \in S. \exists s' \in S . T(s, s')  $)
    \end{itemize}
\end{defframe}

\begin{defframe}{Kripke Structure}{}
    A quadruple $(S, T, I, L)$ is a \textbf{Kripke Structure} if 
    \begin{itemize}
        \item $(S, T, I, L)$ is a finite-state transition system
        \item $L$ is a labelling function $S \to 2^{AP}$(from states to atomic propositions)
    \end{itemize}
\end{defframe}

Atomic propositions are used to represent properties of states, which can be combined them to generate \textit{assertions}.

We can check whether an assertion holds in a given state via the standard $\vDash$ relation - we write $s \vDash F$ if $F$ holds in state $s$:
\begin{itemize}
    \item $s \vDash p \iff p \in L(s)$
    \item $s \not\vDash F \iff p \not\in L(s)$ (!! closed world assumption - anything not explicitly defined as true in $L(s)$ is evaluated as false)
    \item $s \vDash F_{1} \lor F_{2} \iff s \vDash F_{1}\text{ or } s \vDash F_{2}$
    \item $s \vDash F_{1} \land  F_{2} \iff s \vDash F_{1}\text{ and } s \vDash F_{2}$
\end{itemize}

Assertions can be used to express that something bad is not supposed to happen (\textit{safety} properties).

For example, if we had a traffic light system, we could write $\mathtt{assert}(\neg \textcolor{OliveGreen}{\mathtt{green_{1}}} \lor \neg \textcolor{OliveGreen}{\mathtt{green_{2}}})$ to signify that traffic lights \texttt{1} and \texttt{2} cannot be green at the same time.

\section{Explicit-State Model Checking of Assertions}
To algorithmically verify if a system design is correct, we can use \textbf{explicit-state model checking}. This technique systematically enumerates all reachable states of a finite-state transition system, starting from the initial states $I$ and following the transition relation $T$, to verify whether a specified property or assertion $F$ holds in every reachable state. The two primary search strategies used to explicitly traverse and evaluate this state space are Depth-First Search (DFS) and Breadth-First Search (BFS).

\noindent
\begin{minipage}[t]{0.48\textwidth}
\textbf{Depth-First Search (DFS)} \\
\textbf{Algorithm} \textsc{ExplicitReachDFS}
\begin{algorithmic}[1]
\State $open \gets \text{list}(I)$
\State $visited \gets \emptyset$
\While{$open \neq []$}
    \State $s \gets \text{head}(open)$
    \State $open \gets \text{tail}(open)$
    \If{$s \not\models F$} 
        \State \textcolor{red}{error}(s)
    \EndIf
    \For{\textcolor{blue}{each} $s' \in S : (s,s') \in T$}
        \If{$s' \notin visited$}
            \State $visited \gets \{s'\} \cup visited$
            \State $open \gets s' \mathbin{::} open$ \Comment{prepend}
        \EndIf
    \EndFor
\EndWhile
\end{algorithmic}

\vspace{1cm}
\begin{itemize}
    \item keeps only the current stack (\(|\text{open}| < |S|\)): considered more memory-efficient
    \item counterexamples can be quite long
\end{itemize}
\end{minipage}
\hfill
\begin{minipage}[t]{0.48\textwidth}
\textbf{Breadth-First Search (BFS)} \\
\textbf{Algorithm} \textsc{ExplicitReachBFS}
\begin{algorithmic}[1]
\State $open \gets \text{list}(I)$
\State $visited \gets \emptyset$
\While{$open \neq []$}
    \State $s \gets \text{head}(open)$
    \State $open \gets \text{tail}(open)$
    \If{$s \not\models F$} 
        \State \textcolor{red}{error}(s)
    \EndIf
    \For{\textcolor{blue}{each} $s' \in S : (s,s') \in T$}
        \If{$s' \notin visited$}
            \State $visited \gets \{s'\} \cup visited$
            \State $open \gets \text{append}(open, s')$ \Comment{append}
        \EndIf
    \EndFor
\EndWhile
\end{algorithmic}

\vspace{0.5cm}
\begin{itemize}
    \item queue grows very large due to frontier size: consumes significantly more memory
    \item always finds the shortest counterexample
\end{itemize}
\end{minipage}

But assertions only allow us to make sure that nothing bad happens: what if we want to guarantee that something \textit{good} happens?

We need to extend our specification language to consider \textbf{temporal logics}.

\section{Temporal Logics}



\end{document}
